% Claim proof file (REFACTOR TWIN) for row claim_3_7 -- "TwoDisjointNode".
% This is the twin file produced during the `eqViaNodeMap_injective`
% refactor; the existing `claim_3_7_proof_TwoDisjointNode.tex` (itself a
% prior `cdmg_typed_edges` twin) stays untouched until cleanup at
% Phase 7 atomically renames this twin over it.
%
% Mathematically the LN-level proof is unchanged by the refactor: the
% lemma asserts the same triple equality of CDMGs (up to canonical
% carrier bijection) as before, with the same disjointness side
% condition on $W_1, W_2 \ins V$ and the same node-splitting
% construction of def_3_11. The only addition versus the existing
% `cdmg_typed_edges` twin's proof body is a new ``fifth componentwise
% check'' paragraph in each of sub-claims (a) and (b) that justifies
% the strengthened Lean predicate's carrier-map injectivity side
% condition: with the new shape of `eqViaNodeMap` (refactor
% `eqViaNodeMap_injective`) carrying a `Set.InjOn f (G.J \cup G.V)`
% conjunct in addition to the four image-equality conjuncts of the
% prior shape, the proof now also needs to certify that the canonical
% carrier bijection of the statement block is injective when restricted
% to the iterated graph's joint input-output node set. The new
% paragraphs spell this out as a five-cell case analysis on
% $J_{\lp G_{\spl(W_1)} \rp_{\spl(W_2)}} \cup V_{\lp G_{\spl(W_1)} \rp_{\spl(W_2)}}$
% (and its swapped counterpart for sub-claim (b)), using the
% disjointness hypothesis $W_1 \cap W_2 = \emptyset$ to rule out the
% two collision patterns that the LN's identification of carriers
% otherwise leaves implicit. The rest of the proof body --
% statement-restatement block, four componentwise checks on
% $J / V / E / L$, and the closing transitivity argument -- is verbatim
% from the `cdmg_typed_edges` twin.
%
% This file is a sibling of `main.tex` inside `<subsection>/tex/`. The
% `subfiles` package lets it render both standalone and as part of
% `main.tex` via `\subfile{...}`.

\documentclass[main]{subfiles}

\begin{document}
% ref: claim_3_7 (proof, with statement restated for readability)
% title: TwoDisjointNode
\def\rowref{claim\_3\_7}
\phantomsection\label{claim_3_7}

% --- statement (restated) ---------------------------------------------------
% Statement body below is synced from the canonical statement file
% `tex/claim_3_7_statement_TwoDisjointNode.tex`, with one light
% refactor-aware rewording: the L-axiom recitation drops the explicit
% ordered-pair "((v_1, v_2) \in L \Leftrightarrow (v_2, v_1) \in L)"
% symmetric clause and uses the unordered-pair phrasing
% "L consists of unordered pairs of distinct vertices in V"; the
% mathematics is unchanged and the same as in the LN.
%
% NOTE on the iterated-splitting carrier mismatch (Lean encoding).
%
% In a Lean encoding that follows def_3_11 literally and realises the
% tagged copies W^0, W^1 as type-level disjoint sums, the carriers of
% (G_{\spl(W_1)})_{\spl(W_2)} and G_{\spl(W_1 \dcup W_2)} are formally
% distinct types (twice-tagged vs. once-tagged on J \cup V), even though
% the LN identifies them set-theoretically via def_3_11's unsplit-
% injection convention "v^0 := v^1 := v for v \in J \cup (V \sm W)". The
% one-sentence clarification inside the lemma body below ("=" is read up
% to the canonical bijection of carriers) is the load-bearing reading;
% the concrete choice of how to render this identification in Lean
% (CDMG-transport along a node-bijection, CDMG-isomorphism structure,
% or some other equivalence layer) is deferred to the Lean-formalization
% worker. The same wrinkle recurs in later rows of Section 3.2 in which
% iterated node-splitting (or node-splitting composed with hard
% intervention, marginalisation, etc.) is involved, so the encoding
% chosen at the Lean-formalization step here is likely to be re-used.
% See workspace_claim_3_7.md for the encoding-design discussion.

\begin{Lem}[Two disjoint node-splittings commute]
    Let $G = (J, V, E, L)$ be a CDMG (def \ref{def-cdmg}); throughout we use the notation of def \ref{not-cdmg}, recall the typing and axioms of def \ref{def-cdmg} (namely $J \cap V = \emptyset$, $E \ins (J \cup V) \times V$, and $L$ consists of unordered pairs of distinct vertices in $V$), and rely on the node-splitting construction of def \ref{def:G_node-splitting}. Let
    \[ W_1 \ins V \quad\text{and}\quad W_2 \ins V \]
    be two subsets of the output-node set of $G$, subject to the disjointness side condition
    \[ W_1 \cap W_2 \;=\; \emptyset. \]
    All three quantifiers ``for every CDMG $G$'', ``for every $W_1 \ins V$'', and ``for every $W_2 \ins V$ with $W_1 \cap W_2 = \emptyset$'' are read \emph{universally}; in particular the (vacuously disjoint) case $W_1 = W_2 = \emptyset$ is admitted.

    \emph{Reading of ``CDMG'' versus ``CADMG'' in the LN's prose.}
    The LN's literal text contains a typographical inconsistency: it opens with ``Let $G = (J, V, E, L)$ be a CDMG'' but writes ``the CDMG\ldots is the same CADMG\ldots'' in the prose connecting the two iterated splittings. We read this entire row as a statement about \emph{CDMGs} throughout, in line with the addition-to-the-LN for this row and with def \ref{def:G_node-splitting}, whose construction takes a CDMG to a CDMG with no acyclicity attribute claimed (see in particular the closing paragraphs of def \ref{def:G_node-splitting} recording that $G_{\spl(W)}$ can carry $2$-cycles when $G$ has self-loops on $W$, so $G_{\spl(W)}$ is in general not a CADMG even when $G$ is). The asserted equality below is consequently an equality of CDMGs (def \ref{def-cdmg}).

    \emph{Well-typedness of the iterated splitting.}
    The inner two terms of the displayed equality below are iterated node-splittings of the form $\lp G_{\spl(W_i)} \rp_{\spl(W_j)}$ (with $\{i, j\} = \{1, 2\}$), in which $\spl(W_j)$ is applied to the CDMG $G_{\spl(W_i)}$ produced by def \ref{def:G_node-splitting}. For this outer application to be well-typed per def \ref{def:G_node-splitting}'s precondition ``$W \ins V$ of the input CDMG'', we need $W_j$ to be (a subset of) the output-node set of $G_{\spl(W_i)}$, i.e.\
    \[ W_j \;\ins\; V_{\spl(W_i)} \;=\; (V \sm W_i) \;\dcup\; W_i^0 \;\dcup\; W_i^1, \]
    where the equality of carriers is item ii.\ of def \ref{def:G_node-splitting} and $W_i^0 := \lC w^0 \st w \in W_i \rC$, $W_i^1 := \lC w^1 \st w \in W_i \rC$ are the two tagged copies of $W_i$ introduced there. By the disjointness hypothesis $W_1 \cap W_2 = \emptyset$, every $w \in W_j$ satisfies $w \in V$ (since $W_j \ins V$) and $w \notin W_i$ (by disjointness), so $w \in V \sm W_i$. Through def \ref{def:G_node-splitting}'s notational shorthand ``$v^0 := v^1 := v$ for $v \in J \cup (V \sm W_i)$'', the set $V \sm W_i$ is realised as the untagged piece of the disjoint union $V_{\spl(W_i)} = (V \sm W_i) \dcup W_i^0 \dcup W_i^1$, i.e.\ the inclusion
    \[ V \sm W_i \;\hookrightarrow\; V_{\spl(W_i)} \]
    sends $v \mapsto v$ unchanged. Hence
    \[ W_j \;\ins\; V \sm W_i \;\ins\; V_{\spl(W_i)}, \]
    so $\lp G_{\spl(W_i)} \rp_{\spl(W_j)}$ is defined per def \ref{def:G_node-splitting}. The analogous spell-out for the swapped pair (with $i$ and $j$ exchanged) is the same argument with the role of $W_1$ and $W_2$ exchanged, using the same disjointness condition $W_1 \cap W_2 = W_2 \cap W_1 = \emptyset$. The right-hand-side single splitting $G_{\spl(W_1 \dcup W_2)}$ is well-typed per def \ref{def:G_node-splitting} because $W_1 \dcup W_2 \ins V$: by $W_1 \ins V$ and $W_2 \ins V$ together with disjointness, the union $W_1 \cup W_2$ is itself a disjoint union $W_1 \dcup W_2$ and remains a subset of $V$.

    \emph{The conclusion.}
    The CDMGs $\lp G_{\spl(W_1)} \rp_{\spl(W_2)}$, $\lp G_{\spl(W_2)} \rp_{\spl(W_1)}$, and $G_{\spl(W_1 \dcup W_2)}$ (all constructed via def \ref{def:G_node-splitting} as just verified) coincide. Concretely, the LN's displayed triple equality
    \[ \lp G_{\spl(W_1)} \rp_{\spl(W_2)} \;=\; \lp G_{\spl(W_2)} \rp_{\spl(W_1)} \;=\; G_{\spl(W_1 \dcup W_2)} \]
    is the conjunction of the two following equalities of CDMGs:
    \begin{enumerate}[label=(\alph*)]
        \item \emph{Composition with $W_1$ then $W_2$.}
              \[ \lp G_{\spl(W_1)} \rp_{\spl(W_2)} \;=\; G_{\spl(W_1 \dcup W_2)}. \]
        \item \emph{Composition with $W_2$ then $W_1$.}
              \[ \lp G_{\spl(W_2)} \rp_{\spl(W_1)} \;=\; G_{\spl(W_1 \dcup W_2)}. \]
    \end{enumerate}
    Together (a) and (b) yield the LN's ``symmetry under swap'' reading
    \[ \lp G_{\spl(W_1)} \rp_{\spl(W_2)} \;=\; \lp G_{\spl(W_2)} \rp_{\spl(W_1)} \]
    by transitivity of equality on the shared right-hand side $G_{\spl(W_1 \dcup W_2)}$. The equality of CDMGs in (a) and (b) is understood componentwise on the four components $(J, V, E, L)$ of def \ref{def-cdmg}, but the conjunctive unpacking into the four field-by-field equalities is deferred to the proof.

    The equalities in (a) and (b) are equalities of CDMGs (def \ref{def-cdmg}) read \emph{up to the canonical bijection} of the iterated and single-split output-node carriers induced by def \ref{def:G_node-splitting}'s unsplit-injection convention ``$v^0 := v^1 := v$ for $v \in J \cup (V \sm W)$'' (the iterated carrier $V_{\lp G_{\spl(W_1)} \rp_{\spl(W_2)}}$ and the single-step carrier $V_{\spl(W_1 \dcup W_2)}$ are not literally the same set as produced by def \ref{def:G_node-splitting}, but the two are canonically identified by matching corresponding split copies).
\end{Lem}

% --- proof ------------------------------------------------------------------
% Proof body adapted from the LN's \Claude{...} block at
% lecture-notes/lecture_notes/graphs.tex lines 466-492, with four light edits:
% (i) the opening is re-framed to reference the rewrite's (a) and (b) labels
%     (the LN proves only (a) explicitly and folds (b) into a "by symmetry"
%     remark; we keep that structure but make the labelling explicit),
% (ii) the well-typedness of the outer splitting application is cited from
%      the statement block above rather than restated in full (the LN proof
%      spells it out in its first sentence; the rewritten statement spells
%      it out in full, so the proof just refers back),
% (iii) the closing "by symmetry" remark for (b) is expanded into an
%       explicit W_1 <-> W_2 swap + commutativity-of-disjoint-union step,
%       closing with a transitivity-of-equality remark recovering the LN's
%       triple equality through the shared right-hand side, and
% (iv) a new "fifth componentwise check" paragraph -- injectivity of the
%      canonical carrier bijection on the iterated graph's J \cup V -- is
%      inserted at the end of each of sub-claims (a) and (b)'s
%      componentwise verification, just before the "Combining ..." closing
%      lines.  The new paragraph is *not* in the LN's \Claude proof (the
%      LN identifies the iterated and single-step carriers implicitly via
%      the canonical bijection of def_3_11's unsplit-injection convention,
%      so no injectivity certificate is needed in the LN's reading); it is
%      added here to underwrite the strengthened Lean-encoding of the
%      "equality up to canonical bijection of carriers" predicate
%      `eqViaNodeMap` introduced by refactor `eqViaNodeMap_injective`.
%      Mathematically it is a straightforward five-cell case analysis on
%      the iterated graph's J \cup V using the disjointness hypothesis
%      W_1 \cap W_2 = \emptyset; it lives at the LN/Lean interface and is
%      not load-bearing for the LN-level mathematics of the lemma.
\begin{proof}
We prove sub-claim (a), $\lp G_{\spl(W_1)} \rp_{\spl(W_2)} = G_{\spl(W_1 \dcup W_2)}$, by checking the equality of the two CDMGs componentwise on each of the four components $(J, V, E, L)$ of def \ref{def-cdmg} in turn; sub-claim (b), $\lp G_{\spl(W_2)} \rp_{\spl(W_1)} = G_{\spl(W_1 \dcup W_2)}$, is then obtained by an explicit symmetry argument at the end. Both sub-claims rely on the well-typedness check already established in the statement block above to license the outer $\spl(\cdot)$ application: in particular, by the disjointness hypothesis $W_1 \cap W_2 = \emptyset$, every $w \in W_2$ satisfies $w \in V \sm W_1 \ins V_{\spl(W_1)}$ (and symmetrically every $w \in W_1$ satisfies $w \in V \sm W_2 \ins V_{\spl(W_2)}$).

\smallskip\noindent\emph{Notation for split copies.}
For $k \in \{1, 2\}$ and $w \in W_k$, write $w^{0_k}$ and $w^{1_k}$ for the two tagged copies produced by $\spl(W_k)$ per def \ref{def:G_node-splitting}; the subscript $k$ records which of the two splittings introduced the tag. By def \ref{def:G_node-splitting}'s unsplit-injection shorthand, $v^{0_k} := v^{1_k} := v$ for $v \in J \cup (V \sm W_k)$. For an element $v \in J \cup V$ we will also write $v^0$ for the single-step tagging used on $G_{\spl(W_1 \dcup W_2)}$: $v^0 := v^{0_1}$ if $v \in W_1$, $v^0 := v^{0_2}$ if $v \in W_2$, and $v^0 := v$ if $v \in J \cup (V \sm (W_1 \dcup W_2))$ (the three cases are exhaustive and pairwise disjoint by disjointness $W_1 \cap W_2 = \emptyset$); the symmetric convention $v^1$ uses $1_k$ in place of $0_k$. This is the canonical bijection of carriers spelled out at the bottom of the statement block.

\medskip\noindent\textbf{Proof of (a): composition with $W_1$ then $W_2$.}
We verify the equality $\lp G_{\spl(W_1)} \rp_{\spl(W_2)} = G_{\spl(W_1 \dcup W_2)}$ on each of the four components.

\medskip\noindent\emph{Input nodes.}
By item i.\ of def \ref{def:G_node-splitting} applied twice (first to $G$ with $W_1$, then to $G_{\spl(W_1)}$ with $W_2$),
\[
    J_{\lp G_{\spl(W_1)} \rp_{\spl(W_2)}} \;=\; J_{\spl(W_1)} \;=\; J \;=\; J_{\spl(W_1 \dcup W_2)},
\]
where the last equality is item i.\ of def \ref{def:G_node-splitting} applied to the single-step splitting with split set $W_1 \dcup W_2$. \checkmark

\medskip\noindent\emph{Output nodes.}
By item ii.\ of def \ref{def:G_node-splitting} applied to the outer splitting (acting on the input CDMG $G_{\spl(W_1)}$, with split set $W_2 \ins V_{\spl(W_1)}$ as licensed by the well-typedness check),
\[
    V_{\lp G_{\spl(W_1)} \rp_{\spl(W_2)}} \;=\; \lp V_{\spl(W_1)} \sm W_2 \rp \;\dcup\; W_2^{0_2} \;\dcup\; W_2^{1_2}.
\]
Item ii.\ of def \ref{def:G_node-splitting} applied to the inner splitting gives $V_{\spl(W_1)} = (V \sm W_1) \dcup W_1^{0_1} \dcup W_1^{1_1}$. By disjointness $W_1 \cap W_2 = \emptyset$, the elements of $W_2 \ins V$ all lie in $V \sm W_1$, and (via the unsplit-injection identification $V \sm W_1 \hookrightarrow V_{\spl(W_1)}$) they embed into the untagged piece of the above disjoint union, disjointly from $W_1^{0_1} \dcup W_1^{1_1}$. Removing $W_2$ from $V_{\spl(W_1)}$ thus only affects the untagged piece, giving
\[
    V_{\spl(W_1)} \sm W_2 \;=\; (V \sm W_1 \sm W_2) \;\dcup\; W_1^{0_1} \;\dcup\; W_1^{1_1}.
\]
Substituting,
\[
    V_{\lp G_{\spl(W_1)} \rp_{\spl(W_2)}} \;=\; (V \sm W_1 \sm W_2) \;\dcup\; W_1^{0_1} \;\dcup\; W_1^{1_1} \;\dcup\; W_2^{0_2} \;\dcup\; W_2^{1_2}.
\]
On the other hand, item ii.\ of def \ref{def:G_node-splitting} applied to the single-step splitting on $G$ with split set $W_1 \dcup W_2$ gives
\[
    V_{\spl(W_1 \dcup W_2)} \;=\; (V \sm (W_1 \dcup W_2)) \;\dcup\; (W_1 \dcup W_2)^0 \;\dcup\; (W_1 \dcup W_2)^1,
\]
and $V \sm (W_1 \dcup W_2) = V \sm W_1 \sm W_2$ together with the tag-by-tag identification of the single-step copies $(W_1 \dcup W_2)^0 = W_1^{0_1} \dcup W_2^{0_2}$ and $(W_1 \dcup W_2)^1 = W_1^{1_1} \dcup W_2^{1_2}$ (per the canonical bijection of carriers in the notation block above) yields
\[
    V_{\spl(W_1 \dcup W_2)} \;=\; (V \sm W_1 \sm W_2) \;\dcup\; W_1^{0_1} \;\dcup\; W_1^{1_1} \;\dcup\; W_2^{0_2} \;\dcup\; W_2^{1_2}.
\]
The two output-node sets match. \checkmark

\medskip\noindent\emph{Directed edges.}
By item iii.\ of def \ref{def:G_node-splitting} applied to the inner splitting,
\[
    E_{\spl(W_1)} \;=\; \lC v_1^{1_1} \tuh v_2^{0_1} \st v_1 \tuh v_2 \in E \rC \;\cup\; \lC w^{0_1} \tuh w^{1_1} \st w \in W_1 \rC.
\]
Applying the outer $\spl(W_2)$ to the input CDMG $G_{\spl(W_1)}$, item iii.\ of def \ref{def:G_node-splitting} gives
\[
    E_{\lp G_{\spl(W_1)} \rp_{\spl(W_2)}} \;=\; \lC (v_1')^{1_2} \tuh (v_2')^{0_2} \st v_1' \tuh v_2' \in E_{\spl(W_1)} \rC \;\cup\; \lC w^{0_2} \tuh w^{1_2} \st w \in W_2 \rC.
\]
We trace each generator of $E_{\spl(W_1)}$ through the outer tagging.

\smallskip\noindent\emph{Generators from $E$.} For a generator $v_1^{1_1} \tuh v_2^{0_1}$ arising from $v_1 \tuh v_2 \in E$ (so $v_1 \in J \cup V$ and $v_2 \in V$ by def \ref{def-cdmg}), we case-split on each endpoint independently.

For the head $v_1^{1_1}$:
\begin{itemize}
    \item If $v_1 \in W_1$, then $v_1^{1_1} \in W_1^{1_1}$. By the inner-splitting carrier identity $V_{\spl(W_1)} = (V \sm W_1) \dcup W_1^{0_1} \dcup W_1^{1_1}$, the tagged piece $W_1^{1_1}$ is disjoint from the untagged piece $V \sm W_1 \supseteq W_2$, so $v_1^{1_1} \notin W_2$; the outer unsplit-injection $u^{1_2} := u$ for $u \in V_{\spl(W_1)} \sm W_2$ then gives $(v_1^{1_1})^{1_2} = v_1^{1_1} = v_1^1$ (in the $v^1$ notation of the notation block).
    \item If $v_1 \in W_2$, then by disjointness $v_1 \notin W_1$, so the inner unsplit-injection gives $v_1^{1_1} = v_1$; the outer splitting then maps $v_1 \mapsto v_1^{1_2} = v_1^1$.
    \item If $v_1 \in J \cup (V \sm (W_1 \dcup W_2))$, then $v_1 \notin W_1$ and $v_1 \notin W_2$, so $v_1^{1_1} = v_1$ (inner unsplit) and $(v_1)^{1_2} = v_1 = v_1^1$ (outer unsplit).
\end{itemize}
The same case-analysis applied to the tail $v_2^{0_1}$ -- with $0_k$ in place of $1_k$ throughout, and the slightly tighter range $v_2 \in V$ -- yields the tail tagging $v_2^0$ in each of the three cases. Hence the image of $v_1 \tuh v_2 \in E$ under the iterated tagging is exactly $v_1^1 \tuh v_2^0$, and the collection of such images is
\[
    \lC v_1^1 \tuh v_2^0 \st v_1 \tuh v_2 \in E \rC.
\]

\smallskip\noindent\emph{Inner split-arc generators.} For a generator $w^{0_1} \tuh w^{1_1}$ arising from $w \in W_1$, both endpoints lie in $W_1^{0_1} \cup W_1^{1_1}$. By the same disjointness-from-the-untagged-piece argument used in the head case, neither tagged copy lies in $W_2$, so the outer unsplit-injection fixes both: $(w^{0_1})^{0_2} = w^{0_1}$ and $(w^{1_1})^{1_2} = w^{1_1}$. The edge $w^{0_1} \tuh w^{1_1}$ is thus unaffected by the outer splitting. The full set of inner split-arc images is
\[
    \lC w^{0_1} \tuh w^{1_1} \st w \in W_1 \rC \;=\; \lC w^0 \tuh w^1 \st w \in W_1 \rC,
\]
where the rewrite into the unified $w^0$/$w^1$ notation uses the case ``$w \in W_1$'' of the notation block.

\smallskip\noindent\emph{Outer split-arc generators.} The second clause of the outer splitting contributes the fresh split-arc set
\[
    \lC w^{0_2} \tuh w^{1_2} \st w \in W_2 \rC \;=\; \lC w^0 \tuh w^1 \st w \in W_2 \rC,
\]
using the case ``$w \in W_2$'' of the notation block.

\smallskip\noindent\emph{Combination.} Taking the union of the three pieces above,
\begin{align*}
    E_{\lp G_{\spl(W_1)} \rp_{\spl(W_2)}}
        &= \lC v_1^1 \tuh v_2^0 \st v_1 \tuh v_2 \in E \rC \\
        &\quad \cup\; \lC w^0 \tuh w^1 \st w \in W_1 \rC \;\cup\; \lC w^0 \tuh w^1 \st w \in W_2 \rC \\
        &= \lC v_1^1 \tuh v_2^0 \st v_1 \tuh v_2 \in E \rC \;\cup\; \lC w^0 \tuh w^1 \st w \in W_1 \dcup W_2 \rC,
\end{align*}
where the last equality folds the two split-arc sets into a single one indexed by $w \in W_1 \dcup W_2$ (the two index sets are disjoint by hypothesis, so the union is in fact a disjoint union).

On the other hand, item iii.\ of def \ref{def:G_node-splitting} applied to the single-step splitting with split set $W_1 \dcup W_2 \ins V$ yields directly
\[
    E_{\spl(W_1 \dcup W_2)} \;=\; \lC v_1^1 \tuh v_2^0 \st v_1 \tuh v_2 \in E \rC \;\cup\; \lC w^0 \tuh w^1 \st w \in W_1 \dcup W_2 \rC,
\]
under the same $v^0$/$v^1$ tagging conventions. Hence
\[
    E_{\lp G_{\spl(W_1)} \rp_{\spl(W_2)}} \;=\; E_{\spl(W_1 \dcup W_2)}.
\]
\checkmark

\medskip\noindent\emph{Bidirected edges.}
By item iv.\ of def \ref{def:G_node-splitting} applied to the inner splitting,
\[
    L_{\spl(W_1)} \;=\; \lC v_1^{0_1} \huh v_2^{0_1} \st v_1 \huh v_2 \in L \rC.
\]
Applying the outer $\spl(W_2)$, item iv.\ of def \ref{def:G_node-splitting} gives
\[
    L_{\lp G_{\spl(W_1)} \rp_{\spl(W_2)}} \;=\; \lC (v_1')^{0_2} \huh (v_2')^{0_2} \st v_1' \huh v_2' \in L_{\spl(W_1)} \rC.
\]
For a generator $v_1^{0_1} \huh v_2^{0_1}$ (arising from $v_1 \huh v_2 \in L$, with $v_1, v_2 \in V$ by def \ref{def-cdmg}), the same case analysis as for directed edges -- now applied to each $0$-superscripted endpoint independently, instead of the mixed $1_1$-head / $0_1$-tail of the directed-edge case -- yields the endpoint tagging $v_1^0$ and $v_2^0$ in the unified notation of the notation block (per the three cases ``$v \in W_1$'', ``$v \in W_2$'', ``$v \notin W_1 \dcup W_2$''). Hence
\[
    L_{\lp G_{\spl(W_1)} \rp_{\spl(W_2)}} \;=\; \lC v_1^0 \huh v_2^0 \st v_1 \huh v_2 \in L \rC.
\]
On the other hand, item iv.\ of def \ref{def:G_node-splitting} applied to the single-step splitting with split set $W_1 \dcup W_2$ yields exactly the same set under the same $v^0$ tagging:
\[
    L_{\spl(W_1 \dcup W_2)} \;=\; \lC v_1^0 \huh v_2^0 \st v_1 \huh v_2 \in L \rC.
\]
Hence
\[
    L_{\lp G_{\spl(W_1)} \rp_{\spl(W_2)}} \;=\; L_{\spl(W_1 \dcup W_2)}.
\]
\checkmark

\medskip\noindent\emph{Injectivity of the carrier bijection on the iterated graph's node set.}
The canonical bijection of carriers spelled out in the notation block above (and in the closing paragraph of the statement block) sends each element of $J_{\lp G_{\spl(W_1)} \rp_{\spl(W_2)}} \cup V_{\lp G_{\spl(W_1)} \rp_{\spl(W_2)}}$ to its single-step counterpart in $J_{\spl(W_1 \dcup W_2)} \cup V_{\spl(W_1 \dcup W_2)}$ via the rules ``twice-untagged $\mapsto$ untagged'' (i.e., $v \mapsto v$ for $v \in J \cup (V \sm (W_1 \dcup W_2))$) and ``$w^{0_k} \mapsto w^0$, $w^{1_k} \mapsto w^1$'' (for $w \in W_k$, $k \in \{1, 2\}$). We verify that this bijection is moreover \emph{injective} on the source side -- i.e., distinct elements of the iterated graph's joint input-output node set are sent to distinct elements of the single-step graph's joint input-output node set -- by a case analysis on the five-cell partition of the source.

Combining the input-node and output-node identities verified above with item i.\ of def \ref{def:G_node-splitting}'s ``$J_{\spl(\cdot)} = J$'' identity for both the inner and outer splittings, the joint input-output set of the iterated graph decomposes as the disjoint union
\[
    J_{\lp G_{\spl(W_1)} \rp_{\spl(W_2)}} \cup V_{\lp G_{\spl(W_1)} \rp_{\spl(W_2)}}
    \;=\; \lp J \cup (V \sm W_1 \sm W_2) \rp \;\dcup\; W_1^{0_1} \;\dcup\; W_1^{1_1} \;\dcup\; W_2^{0_2} \;\dcup\; W_2^{1_2},
\]
where the first cell collects every element that is untagged at both splittings (whose single-step image is itself untagged, written $v$), and the remaining four cells are the tagged copies introduced by the two splittings (with single-step images $w^0$ resp.\ $w^1$ in $(W_1 \dcup W_2)^0$ resp.\ $(W_1 \dcup W_2)^1$ per the notation block). Disjointness of the five cells inside the iterated graph's joint input-output set follows from items i.\ and ii.\ of def \ref{def:G_node-splitting} together with the disjointness hypothesis $W_1 \cap W_2 = \emptyset$ (the disjointness side condition on $J \cap V = \emptyset$ from def \ref{def-cdmg} disposes of the $J$-vs.-$V$ piece of the first cell's pairwise disjointness with the four tagged cells, and the tagged copies are pairwise disjoint from the untagged piece by item ii.\ of def \ref{def:G_node-splitting}).

\emph{Within-cell injectivity.} On each of the five cells the bijection acts by a single composed-constructor chain that is injective in the underlying $w$- or $v$-label: ``$v \mapsto v$'' on $J \cup (V \sm W_1 \sm W_2)$, ``$w \mapsto w^0$'' on $W_1^{0_1}$ and on $W_2^{0_2}$, ``$w \mapsto w^1$'' on $W_1^{1_1}$ and on $W_2^{1_2}$. Each of these underlying maps is injective in $w$ (resp.\ $v$) by item ii.\ of def \ref{def:G_node-splitting}'s disjoint-union typing of $V_{\spl(W_1 \dcup W_2)}$ (the single-step copies $w^0$, $w^1$ are pairwise distinct as $w$ ranges over $W_1 \dcup W_2$, by construction of the disjoint union).

\emph{Across-cell injectivity.} It remains to show that no collision occurs between two elements drawn from two different cells of the five-cell partition. The bijection's image on the first cell lands in the untagged piece $J \cup (V \sm (W_1 \dcup W_2))$ of the single-step graph's joint input-output set, while its image on the four tagged cells lands in $(W_1 \dcup W_2)^0 \cup (W_1 \dcup W_2)^1$; these two pieces are disjoint by item ii.\ of def \ref{def:G_node-splitting} (the untagged piece is exactly the complement of the two tagged copies in $V_{\spl(W_1 \dcup W_2)}$, and they are joined disjointly with $J$), so any cross-collision involving the first cell is impossible. Among the four tagged cells, the structural cross-collision patterns to rule out are:
\begin{itemize}
    \item \emph{$W_1^{0_1}$ vs.\ $W_2^{0_2}$.} Iterated-graph images $w_1^{0_1}$ (for some $w_1 \in W_1$) and $w_2^{0_2}$ (for some $w_2 \in W_2$) map under the canonical bijection to $w_1^0$ and $w_2^0$, both inside $(W_1 \dcup W_2)^0$. An equality $w_1^0 = w_2^0$ in $(W_1 \dcup W_2)^0$ forces $w_1 = w_2$ by injectivity of the single-step $w \mapsto w^0$ tagging (item ii.\ of def \ref{def:G_node-splitting}, the disjoint-union construction of $(W_1 \dcup W_2)^0$ from $\{w \st w \in W_1 \dcup W_2\}$). But $w_1 \in W_1$ and $w_2 \in W_2$ with $W_1 \cap W_2 = \emptyset$ rules this equality out -- contradiction. So no $W_1^{0_1}$-vs.-$W_2^{0_2}$ collision occurs. \emph{This is the load-bearing use of the disjointness hypothesis $W_1 \cap W_2 = \emptyset$.}
    \item \emph{$W_1^{1_1}$ vs.\ $W_2^{1_2}$.} Same argument with the $0$-tag replaced by the $1$-tag throughout (using the $w \mapsto w^1$ injectivity of item ii.\ on the $(W_1 \dcup W_2)^1$ side, and the same disjointness hypothesis $W_1 \cap W_2 = \emptyset$).
    \item \emph{$^0$-tagged cell vs.\ $^1$-tagged cell ($W_1^{0_1}$ vs.\ $W_1^{1_1}$, $W_1^{0_1}$ vs.\ $W_2^{1_2}$, $W_2^{0_2}$ vs.\ $W_1^{1_1}$, $W_2^{0_2}$ vs.\ $W_2^{1_2}$).} On the source side, these are pairs of cells with one mapping into $(W_1 \dcup W_2)^0$ and the other into $(W_1 \dcup W_2)^1$. By item ii.\ of def \ref{def:G_node-splitting}, the two single-step tagged copies $(W_1 \dcup W_2)^0$ and $(W_1 \dcup W_2)^1$ are disjoint by construction (the two-copies clause of the disjoint union ``$(V \sm W) \dcup W^0 \dcup W^1$''), so no $^0$-image can equal any $^1$-image. No collision in any of these four sub-patterns.
\end{itemize}
All cross-cell collision patterns have been ruled out, so the canonical bijection of carriers is injective on $J_{\lp G_{\spl(W_1)} \rp_{\spl(W_2)}} \cup V_{\lp G_{\spl(W_1)} \rp_{\spl(W_2)}}$. \checkmark

\smallskip
Combining the four componentwise equalities above with the carrier-map injectivity check just established, the CDMGs $\lp G_{\spl(W_1)} \rp_{\spl(W_2)}$ and $G_{\spl(W_1 \dcup W_2)}$ agree field by field on each of $(J, V, E, L)$ -- with the four image-equality identities understood up to the canonical bijection of carriers, which is in turn injective on the iterated graph's joint input-output set -- and are therefore equal as CDMGs (def \ref{def-cdmg}). This establishes sub-claim (a).

\medskip\noindent\textbf{Proof of (b): composition with $W_2$ then $W_1$.}
Sub-claim (b) is obtained by swapping the roles of $W_1$ and $W_2$ in the argument for (a). The five-component verification above (four image-equality identities plus the carrier-bijection injectivity check) depends on the pair of subsets only through their symbolic roles: the input-node identity holds because each splitting leaves $J$ unchanged, regardless of the split set; the output-node identity, the directed-edge identity, and the bidirected-edge identity each follow from items ii.--iv.\ of def \ref{def:G_node-splitting} together with the disjointness hypothesis $W_1 \cap W_2 = \emptyset$, which is symmetric in $(W_1, W_2)$ ($W_1 \cap W_2 = \emptyset \Leftrightarrow W_2 \cap W_1 = \emptyset$); and the carrier-bijection injectivity check follows from the same five-cell case analysis with the $^0$-vs.-$^0$ and $^1$-vs.-$^1$ collision patterns ruled out by the same disjointness hypothesis (also symmetric in $(W_1, W_2)$). The well-typedness side-condition needed to license the outer $\spl(\cdot)$ application also has a symmetric counterpart, established in the statement block: every $w \in W_1$ satisfies $w \in V \sm W_2 \ins V_{\spl(W_2)}$.

\medskip\noindent\emph{Injectivity of the carrier bijection on the iterated graph's node set (symmetric counterpart).}
For sub-claim (b)'s iterated graph $\lp G_{\spl(W_2)} \rp_{\spl(W_1)}$, the carrier-bijection injectivity check is the same as in (a) with the substitution $W_1 \leftrightarrow W_2$ throughout: the joint input-output set of the iterated graph decomposes as the five-cell disjoint union
\[
    J_{\lp G_{\spl(W_2)} \rp_{\spl(W_1)}} \cup V_{\lp G_{\spl(W_2)} \rp_{\spl(W_1)}}
    \;=\; \lp J \cup (V \sm W_1 \sm W_2) \rp \;\dcup\; W_2^{0_2} \;\dcup\; W_2^{1_2} \;\dcup\; W_1^{0_1} \;\dcup\; W_1^{1_1},
\]
with the first two tagged cells $W_2^{0_2}, W_2^{1_2}$ now realised by the (inner) $\spl(W_2)$ application and the last two $W_1^{0_1}, W_1^{1_1}$ by the (outer) $\spl(W_1)$ application -- the underlying \emph{set} of the five cells is the same as in (a), but their realisation as iterated-tagged copies inside the iterated graph swaps the inner/outer roles. The canonical bijection still sends each tagged cell $W_k^{\epsilon_k}$ to its single-step image $W_k^{\epsilon}$ inside $(W_1 \dcup W_2)^{\epsilon}$ ($\epsilon \in \{0, 1\}$, $k \in \{1, 2\}$), and the untagged cell to itself. The within-cell injectivity argument is unchanged, and the across-cell analysis uses the same disjointness hypothesis $W_2 \cap W_1 = W_1 \cap W_2 = \emptyset$ to rule out the $^0$-vs.-$^0$ collision pattern ($W_2^{0_2}$ vs.\ $W_1^{0_1}$, which would force a common $w \in W_1 \cap W_2$) and the $^1$-vs.-$^1$ collision pattern ($W_2^{1_2}$ vs.\ $W_1^{1_1}$, same argument), while the remaining $^0$-vs.-$^1$ cross-cell pairs and the untagged-cell-vs.-tagged-cell pairs are ruled out by the same item-ii.\ disjoint-pieces argument as in (a). Hence the canonical bijection of carriers is injective on $J_{\lp G_{\spl(W_2)} \rp_{\spl(W_1)}} \cup V_{\lp G_{\spl(W_2)} \rp_{\spl(W_1)}}$. \checkmark

\medskip
Re-running the five-component check above (four image-equality identities plus the carrier-bijection injectivity check) with the substitution $W_1 \leftrightarrow W_2$ throughout therefore yields
\[
    \lp G_{\spl(W_2)} \rp_{\spl(W_1)} \;=\; G_{\spl(W_2 \dcup W_1)}.
\]
This is precisely sub-claim (a) instantiated at the swapped pair $(W_2, W_1)$. We emphasise that sub-claim (b) is \emph{not} a literal instance of (a): the iterated splitting on the left-hand side of (b) is built from the swapped pair $(W_2, W_1)$, and the right-hand-side disjoint union appearing in (a) for that swapped pair is $W_2 \dcup W_1$, not the $W_1 \dcup W_2$ in the statement of (b). The two right-hand sides agree by commutativity of disjoint union, $W_2 \dcup W_1 = W_1 \dcup W_2$, which gives
\[
    G_{\spl(W_2 \dcup W_1)} \;=\; G_{\spl(W_1 \dcup W_2)}
\]
as CDMGs: def \ref{def:G_node-splitting} depends on $W$ only through the underlying set $W \ins V$ (each of items i.--iv.\ refers to $W$ only via the operations $V \sm W$, $W^0$, $W^1$, and the membership predicate ``$\cdot \in W$''), and equal sets produce equal CDMGs componentwise. Chaining the two equalities,
\[
    \lp G_{\spl(W_2)} \rp_{\spl(W_1)} \;=\; G_{\spl(W_2 \dcup W_1)} \;=\; G_{\spl(W_1 \dcup W_2)},
\]
establishes sub-claim (b).

\medskip\noindent\emph{Recovery of the LN's triple equality.}
By sub-claims (a) and (b) -- both of which equate their respective iterated splitting with the single right-hand side $G_{\spl(W_1 \dcup W_2)}$ -- transitivity of equality on that shared right-hand side yields
\[
    \lp G_{\spl(W_1)} \rp_{\spl(W_2)} \;=\; G_{\spl(W_1 \dcup W_2)} \;=\; \lp G_{\spl(W_2)} \rp_{\spl(W_1)},
\]
which is the LN's stated triple equality and in particular recovers the swap-symmetry reading $\lp G_{\spl(W_1)} \rp_{\spl(W_2)} = \lp G_{\spl(W_2)} \rp_{\spl(W_1)}$.
\end{proof}

\end{document}
